% 图 3.2 —— **结构化重画**（2026-09-25，替换掉 `checks/emit_hand.py` 的逐段拟合稿）
%%
% 结构：两条虚线渐近线 + 一条实线 y 轴 + 一条横轴 + 一条光滑曲线 y = x/(1−x²) + 1 个公式标签。2026-09-25：虚线各段线宽由 2.0/3.0/4.0/4.5 混用统一为 4.0； y 轴由 2 段合并为 1 条
%%
% 旧版是**描摹**：把每一段笔画单独拟合后发出来，连「箭头头」都被当成"一条很粗的短线"描
%   （图 81.5 上一度出现 `\draw[line width=7.0pt] (586,277)--(586,270);` 这种 7px 长的疙瘩），
%   渲染出来是黑疙瘩而不是箭头。本版改成**看懂结构后用 TikZ 画**。
%%
% 做法、约定、已踩过的坑 → `＜工作区＞\checks\redraw\README_重画流程.md`
% 本图圆/椭圆的中心与半径、标签的文字与位置，沿用 probe 的脚手架（那部分它是对的）。
%%
% ⚠ 本图 tikzpicture 是 y=-1pt（**y 翻转**）⇒ `arc` 的角向"递减"才是屏幕逆时针；
%   同一根线上的多段弧必须同向。改完务必拿原书并排、放大核箭头朝向。
%%
\documentclass[12pt]{standalone}
\usepackage{fontspec}
\setsansfont{Arial}
\newfontfamily\mfallback{DejaVu Sans}
\usepackage{tikz}
\begin{document}
\begin{tikzpicture}[x=1pt, y=-1pt, line width=4.0pt, line cap=round, line join=round]
  % 2026-09-25：两条虚线改为**一条 dash pattern**（原为 20 段硬编码直线，
  %   段长 27.7、间隔 13.6 均由原书实测；原各段线宽在 2.0/3.0/4.0/4.5 间乱跳，现统一 4.0）。
  \draw[line width=4.0pt, dash pattern=on 27.7pt off 13.6pt] (80,12) -- (80,731);
  \draw[line width=4.0pt, dash pattern=on 27.7pt off 13.6pt] (457,8) -- (457,730);
  %% 2026-09-25：虚线统一线宽 4.0（原为 2.0/3.0/4.0/4.5 混用）；y 轴两段合一
  \draw[line width=4.0pt] (269.5,21.0) -- (269.5,738.0);

  %% ════ 直线段（prims2 的 strokes）════
  \draw[line width=4.0pt] (270.0,377.0) -- (270.0,21.0);
  \draw[line width=6.0pt] (97.0,739.0) -- (100.0,727.2);
  \draw[line width=4.0pt] (269.0,381.0) -- (268.0,738.0);

  %% ════ 横轴（合并为一条，宽度统一 4.0）════
  %% 旧稿把它拆成 6 段，且原点附近两段被量成 8.0pt（曲线与轴重叠处墨迹变厚所致），
  %% 成品上表现为轴中间有个鼓包。实测轴厚各处 3.4~4.5，取 4.0。
  %% 轴在原书里略有倾斜（左端 y≈380.5、右端 y≈377.5），此处照抄。
  \draw[line width=4.0pt] (16.0,380.6) -- (529.0,377.6);

  %% ════ 曲线 y = x/(1 − x²) ════
  %% 旧稿是「3 段贝塞尔 + 4 段直线」拼的：近原点处是直线段、线宽还在 6/8/5.5 之间乱跳，
  %% 所以看着不光滑。现改为**一条平滑曲线**（61 个等弧长点 + plot[smooth]）。
  %%
  %% 坐标怎么来的（走了两个弯路，记下来）：
  %%   ① 「逐列取中心线」在**近竖直段失效** —— 同一列里 y 跨几十上百单位，均值是垃圾，
  %%      曲线两端的陡峭部分会被抹成斜的。所以改用**骨架化**（skimage.skeletonize）。
  %%   ② 骨架会被排除带切断，要**分段走线再按端点邻近拼接**（双向贪心）。
  %%   ③ 排除带的位置必须先量准：三条竖线实测在 x≈76.6 / 266.4 / 453.0（且各自略有倾斜），
  %%      横轴从 (16,380.6) 斜到 (529,377.6) —— 原稿用的 80/269.5/457.5 偏了 3 单位。
  %% 最终这条链：弧长 847 单位，x 从 95.8 到 431.6；交点处斜率最小（−0.46）向两侧单调增大
  %% （−22.0 → −0.46 → −17.5），无拐点。
  \draw[line width=6.0pt, line cap=round]
    plot[smooth] coordinates {
(95.8,738.6) (96.4,724.4) (97.1,710.2) (97.8,696.0) (98.5,681.9)
      (99.4,667.8) (100.5,653.7) (101.8,639.7) (103.4,625.7) (105.2,611.8)
      (107.3,598.0) (109.9,584.2) (112.8,570.4) (116.2,556.8) (120.1,543.2)
      (124.6,529.8) (129.6,516.4) (135.2,503.3) (141.4,490.4) (148.3,477.9)
      (155.9,465.8) (164.1,454.3) (173.1,443.4) (182.7,433.3) (193.2,423.9)
      (204.3,415.3) (215.9,407.4) (227.9,400.0) (240.3,393.2) (252.8,386.8)
      (265.5,380.7) (278.1,374.8) (290.7,368.9) (303.0,362.6) (314.9,355.6)
      (326.5,347.7) (337.5,338.9) (347.8,329.2) (357.4,318.8) (366.2,307.7)
      (374.2,295.9) (381.4,283.7) (387.9,271.0) (393.7,258.0) (398.9,244.7)
      (403.6,231.3) (407.7,217.7) (411.3,204.1) (414.5,190.4) (417.2,176.6)
      (419.6,162.8) (421.7,148.9) (423.4,134.9) (425.0,120.9) (426.2,106.9)
      (427.4,92.8) (428.3,78.7) (429.2,64.6) (430.0,50.5) (430.8,36.3)
      (431.6,22.2)
    };

  %% ════ 标签 ════
  %% 原书这里是**一条公式** `y = x/(1 − x²)`（第 43 页，实测墨迹框 39.31×6.69 pt）。
  %% 旧稿（emit_hand.py 自动稿）把它拆成 8 个独立字形，并且认错三处：
  %%   · `=` 被当成减号（还额外抽出一条 504..517 的线段图元）
  %%   · `/` 被认成斜体 `f`
  %%   · 上标 `²` 整个丢掉
  %% 结果读起来是 `y −x f(1 − x')` —— 内容错的，不是格式错的。现改为**单节点整条公式**。
  %% 字号与锚点由「编译→量→修正」收敛到原书实测框；画布换算（用坐标轴标定：
  %%   实线 y 轴 canvas x=269.5 ↔ 原书 152.2pt、虚线渐近线 canvas 80/457.5 ↔ 原书 110.4/193.0pt，
  %%   三处独立吻合 ⇒ canvas = (book − 96.4) × 4.541 + 16）
  \node[anchor=west,inner sep=0pt,fill=white,font=\fontsize{30.0}{37.5}\selectfont]
    at (476.5,198.9) {\textsf{\textbf{\textit{y} = \textit{x}/(1 − \textit{x}\textsuperscript{2})}}};


  % 钉死画布：否则超出的图元/控制点会把页面撑大，1:1 回比就失准
  \pgfresetboundingbox
  \path[use as bounding box] (0,0) rectangle (672,754);
\end{tikzpicture}
\end{document}

% ⚠ 待核清单（probe 拿不准，人眼过一遍）：
%   · 未识别/跳过：⑤ 跳过/未识别的字形
